\documentclass[11pt]{article}
\usepackage{fontspec,amsmath,amssymb,amsthm,hyperref}
\usepackage[margin=1in]{geometry}
\setlength{\emergencystretch}{3em}
\title{Project transport convergence into parameter RMS error}
\author{LangevinLean mathematical exposition}
\date{}
\begin{document}\maketitle
\noindent Source module: \texttt{UnadjustedLangevin/Projection.lean}. Compile this proof document with XeLaTeX or LuaLaTeX. The accompanying HTML places the complete Lean declarations beside every proof.


For a continuous linear contraction \(T:H\to K\), reference point \(b\), and equilibrium projected radius \(c\), the parameter error obeys \((\int\|Tx-b\|^2d\mu)^{1/2}\le W_2(\mu,\pi)+c\). This derives the estimate directly from couplings.


\section{Minkowski on an arbitrary coupling}


Define \(R(\mu)=\|Tx-b\|_{L^2(\mu)}\) as an extended norm, and its real-valued version by conversion. The operator bound \(\|T\|\le1\) gives \(\|T(X-Y)\|_2\le\|X-Y\|_2\). Decompose \(TX-b=T(X-Y)+(TY-b)\), apply Minkowski under each coupling, and use the marginal identities to obtain \(R(\mu)\le\operatorname{cost}(\gamma)+R(\pi)\). Infimizing over the nonempty coupling family gives \(R(\mu)\le W_2(\mu,\pi)+R(\pi)\).



\noindent Lean declarations: \texttt{projectedRMS}, \texttt{projectedRMSReal}, \texttt{projectedDifference\_le}, \texttt{projectedRMS\_le\_coupling}, \texttt{projectedRMS\_le\_wasserstein2}.

\section{Finiteness and the second-moment interpretation}


If the transport distance is finite and \(R(\pi)\le\operatorname{ofReal}(c)\), the preceding upper bound proves \(R(\mu)<\infty\). With this finite-norm premise, the \(L^2\) norm formula gives \(R_{\mathbb R}(\mu)=\sqrt{\int\|Tx-b\|^2d\mu}\). In particular, explicit \(L^2\) membership and the second-moment bound \(\int\|Tx-b\|^2d\pi\le c^2\), \(c\ge0\), imply \(R(\pi)\le\operatorname{ofReal}(c)\).



\noindent Lean declarations: \texttt{projectedRMS\_ne\_top}, \texttt{projectedRMSReal\_eq\_sqrt\_integral}, \texttt{projectedRMS\_le\_of\_secondMoment}.

\section{Convert the finite bound to real quantities}


Selecting a finite coupling ensures that the infimum family is nonempty. The extended inequality \(R(\mu)\le W_2(\mu,\pi)+\operatorname{ofReal}(c)\) has a finite right side. Monotonicity of conversion to \(\mathbb R\), together with \(c\ge0\), yields \(R_{\mathbb R}(\mu)\le W_{2,\mathbb R}(\mu,\pi)+c\). Thus neither the real integral nor the real-valued distance can mask an infinite quantity.



\noindent Lean declarations: \texttt{projectedRMSReal\_le\_wasserstein2Real\_add}.

\end{document}
